\subsection{Особенности оптимизации 3DGS}

Из прошлого раздела были получены параметры, однозначно описывающие гауссианы, и при этом удобные для оптимизации:
%
\[
	\Theta = (\mu_i, s_i, q_i, \alpha_i, \mathrm{sh}_i),
\]
%
которые вместе со всеми предыдущими выкладками могут быть непосредственно использованы для вычисления итогового цвета пикселя изображения с камеры в точке $\mathrm o$ (которой соответствует матрица переноса $-C$) и поворотом $W^{-1}$. Для начала изображение разбивается на блоки $16\times16$ и для каждой гауссианы, попавшей в этот блок происходит сортировка, чтобы правильно произвести alpha blending, а затем применяется формула:
%
\[
	\hat I(u, v) = \sum_i T_i C_i G'_i(u, v),
\]
\[
	T_i = \alpha_i \prod_j (1 - \alpha_j) \text{ по всем перекрывающим гауссианам}
\]
\[
	C_i = \sum_{l=0}^N \sum_{m=-l}^l \mathrm{sh}_i[l][m] \cdot Y_{lm}(\mu_i - \mathrm o),
\]
\[
	G'_i(u, v) = e^{-\frac12 ( (u, v) - \mu_i'^t ) \Sigma'^{-1} ( (u, v)^t - \mu_i' )},
\]
\[
	\mu_i' = P( W\mu_i + C ), \quad
	\Sigma_i' = PJW \Sigma_i (PJW)^t, \quad
	\Sigma_i = R(q_i)S(s) S^t(s)R(q_i)^t
\]
%
где матрицы $P, J, W, R, S$ были описаны в предыдущем разделе.

Теперь по этой формуле можно вычислить и градиент по каждому из параметров. Применяются следующие техники оптимизации.

\textbf{Последовательное усложнение сферических гармоник.} В начале обучения считается (и оптимизируется) только сферические гармоники 0 степени (константа), затем, через определенные промежутки шагов добавляется по степени.

\textbf{Прунинг.} По заданным границам, если непрозрачность гауссианы $\alpha$ маленькая, центр $\mu$ сильно удален или размер $\Sigma$ превосходит допустимый, такая гауссиана удаляется, так как считается бесполезной.

\textbf{Уплотнение.} Гауссианы, для которых градиент по смещению превышает определенный порог, происходит одно из двух: если гауссиана небольшая, то она просто клонируется и смещается в сторону градиента, а если большая, то разделяется на 2 (с уменьшением размера), а центры двух полученных гауссиан размещаются по закону плотности исходной. Такой подход позволяет увеличивать детализацию там, где она требуется.

Этот подход предлагался в оригинальной статье, хотя в этой области ведуться исследования. Одной из удачных эвристик стоит упомянуть работу \cite{kheradmand20253dgaussiansplattingmarkov}.